\section{Model}\label{sec:ModelBasic}

The model consists of heterogeneous agents featuring durable and nondurable goods. Its objective is to provide a transparent benchmark durable--nondurable problem that nests as a special case of the richer environment with search, collateralized borrowing, maintenance costs, and permanent income risk.

In short, the timeline of the model is as follows. The agent starts each period with beginning-of-period durable stock ($n_t$) and market resources ($m_t$). The agent solves the adjuster and keeper problems separately and chooses optimal levels of nondurable consumption $c_t$, durable stock $d_t$, and assets $a_t$. The discrete keep-versus-adjust decision is then resolved by a hard maximum over the two conditional value functions. At the end of the period, the agent holds a durable stock $d_t$ and liquid assets $a_t$. Between period $t$ and $t+1$, the transitory income shock materializes, interest is received, and the durable stock depreciates.

Each household is characterized by two state variables: beginning-of-period market resources $m_t$ and the stock of durables $n_t$. Following \cite{Carroll1992-bh}, all variables are normalized by permanent income and written in lowercase letters. There are no permanent income shocks and no permanent-income growth. Preferences exhibit constant relative risk aversion and combine nondurable consumption $c_t$ and the service flow of durables $d_t$ using a Cobb-Douglas aggregator\footnote{Using this form is standard (eg see \cite{Berger2015-kp, McKay2021-ib}) which is in line with micro-evidence \cite{Berger2018-yb}. An alternative approach is to use a CES aggregator as in \cite{Mankiw1983-za}}:
\begin{align}
	u(c_t, d_t) &= \frac{\big(c_t^\alpha d_t^{1-\alpha}\big)^{1-\rho}}{1-\rho}, \qquad \alpha \in (0,1), \ \rho > 0.
\end{align}

At the start of the period, the household observes $(m_t, n_t)$ and decides whether to keep the durable stock unchanged or to adjust. If no adjustment occurs, the household sets $d_t = n_t$ and allocates $m_t$ between nondurable consumption and savings. If the household adjusts, it sells the entire beginning-of-period stock at no price risk but incurs a proportional adjustment cost $\tau \in (0,1)$. The resulting liquidity is
\begin{align}
	\ell_t = m_t + (1-\tau)n_t,
\end{align}
which can then be allocated across nondurable consumption, a new durable stock $d_t$, and savings $a_t$.

The household's problem is the hard discrete choice between keeping and adjusting:
\begin{align}
	v_t(m_t,n_t) = \max\Big\{ v_t^{\text{keep}}(m_t,n_t),\; v_t^{\text{adj}}(\ell_t)\Big\},
\end{align}
where $v_t^{\text{keep}}$ and $v_t^{\text{adj}}$ are the value functions conditional on keeping or adjusting.

Labor income is idiosyncratic and purely transitory. Normalized income is
\begin{align}
	y_t = \theta_t,
\end{align}
where $\theta_t$ is i.i.d.\ log-normally distributed with $\mathbb{E}[\theta_t] = 1$ and a strictly positive standard deviation. There is no unemployment risk.

Assets earn a single interest rate $r$, so the gross interest factor is $\mathcal{R} = 1+r$. The durable stock depreciates at rate $\delta \in (0,1)$. There are no maintenance costs and no operation costs. Therefore, next-period states evolve according to
\begin{align}
	m_{t+1} &= \mathcal{R} a_t + \theta_{t+1},\\
	n_{t+1} &= (1-\delta)\, d_t.
\end{align}

Borrowing is prohibited: the loan-to-value ratio is zero and the lower bound on assets is $\underline{a}=0$, so end-of-period assets must satisfy $a_t \ge 0$.

Hence, the Bellman equations conditional on keep and adjust are:
\begin{align}
	v_t^{\text{keep}}(m_t,n_t) &= \max_{c_t} \ u(c_t ,d_t) + \beta \mathbb{E}_t \big[ v_{t+1}(m_{t+1},n_{t+1}) \big],\\
	& \text{s.t. } a_t = m_t - c_t \geq 0,\\
	& \qquad d_t = n_t,\\
	& \qquad m_{t+1} = \mathcal{R} a_t + \theta_{t+1},\\
	& \qquad n_{t+1} = (1-\delta)\, d_t.\\[1em]
	v_t^{\text{adj}}(\ell_t) &= \max_{c_t,d_t} \ u(c_t,d_t) + \beta \mathbb{E}_t \big[ v_{t+1}(m_{t+1},n_{t+1}) \big],\\
	& \text{s.t. } a_t = \ell_t - c_t - d_t \geq 0,\\
	& \qquad m_{t+1} = \mathcal{R} a_t + \theta_{t+1},\\
	& \qquad n_{t+1} = (1-\delta)\, d_t.
\end{align}

With the following first-order conditions with respect to nondurable consumption and durable stock:
\begin{align}
	\frac{\partial v_t^{adj}}{\partial c_t} &= u_c(c_t, d_t) \geq \beta \mathcal{R}\, \mathbb{E}_t \big[v_{t+1}^m\big], \label{eq:FOC_c_basic} \\
	\frac{\partial v_t^{adj}}{\partial d_t} &= u_d(c_t, d_t) \geq \beta \mathcal{R}\, \mathbb{E}_t\big[v_{t+1}^m\big] - \beta (1 - \delta)\, \mathbb{E}_t\big[v_{t+1}^n\big], \label{eq:FOC_d_basic}
\end{align}

Equations~\ref{eq:FOC_c_basic} and~\ref{eq:FOC_d_basic} are stated as KKT conditions. Equality holds when the borrowing constraint does not bind, and the left-hand side strictly exceeds the right-hand side when it does. Equation~\ref{eq:FOC_c_basic} equates the marginal utility of nondurable consumption with the discounted shadow value of savings. Equation~\ref{eq:FOC_d_basic} shows that the marginal utility of durables equals the net user cost: the first term on the right is the discounted opportunity cost $\mathcal{R}$ of diverting a unit of resources into durables rather than liquid assets, and the second term is the discounted benefit of carrying the remaining durable stock $(1-\delta)d_t$ into the next period.

I describe the solution method in appendix \ref{sec:appdx_model_basic}. In short, I solve the keeper problem using EGM with an upper envelope algorithm and the adjuster using root-finding and FUES to select the optimal points. The combined policy is the argmax of the keep and adjust branches. Furthermore, I calculate end-of-period variables using continuous shock points for the transitory income distribution instead of approximating it by only a few points.

\subsection{Extension: Taste Shocks}\label{subsec:ModelBasicTaste}

The hard maximum in the discrete keep-versus-adjust decision produces a kink in the value function and a discontinuity in the policy functions at the state combinations $(m_t,n_t)$ where the household is indifferent between the two branches. Following \cite{Khan2019-ew}, this subsection smooths that discontinuity by introducing an idiosyncratic, log-normally distributed, multiplicative taste shock on the adjustment margin, while leaving search frictions aside (these are introduced separately in Section~\ref{subsec:ModelBasicSearch}). Unlike the Extreme Value Type I taste shocks commonly used elsewhere in the discrete-continuous choice literature to randomize the discrete outcome directly, this log-normal shock can take a continuum of values and, once search is added in Section~\ref{subsec:ModelBasicSearch}, feeds into the agent's optimal search-effort decision rather than mechanically randomizing the choice.

At the start of period $t$, after the transitory income shock $\theta_t$ has materialized, the household draws an i.i.d.\ taste shock $\xi_t$ that is specific to the value of adjusting. The shock is log-normally distributed with mean one,
\begin{align}
	\log \xi_t \sim \mathcal{N}\Big(-\frac{\sigma_\xi^2}{2},\, \sigma_\xi^2\Big), \qquad \mathbb{E}[\xi_t] = 1, \label{eq:taste_dist}
\end{align}
independent of $\theta_t$ and independent across periods. The shock can be interpreted as an idiosyncratic realization of the enjoyment (or hassle) of shopping for a new durable in the current period: an unusually good draw makes adjusting relatively more attractive, and an unusually poor draw makes it relatively less attractive, regardless of the underlying state $(m_t,n_t)$.

The conditional value functions $v_t^{\text{keep}}(m_t,n_t)$ and $v_t^{\text{adj}}(\ell_t)$ retain exactly the same Bellman definitions as in the baseline model above. Because $\xi_t$ multiplies the entire conditional value from the perspective of the discrete comparison but does not enter the inner maximization over $c_t$ (and $d_t$ for the adjuster), it does not affect the first-order conditions~\eqref{eq:FOC_c_basic}--\eqref{eq:FOC_d_basic}: a strictly positive scalar multiplying the adjuster's objective leaves its argmax, $c_t^{\text{adj}}$ and $d_t^{\text{adj}}$, unchanged. Instead, $\xi_t$ enters only the discrete comparison itself,
\begin{align}
	v_t(m_t,n_t,\xi_t) = \max\Big\{ v_t^{\text{keep}}(m_t,n_t),\; \xi_t\, v_t^{\text{adj}}(\ell_t) \Big\}. \label{eq:taste_discrete}
\end{align}
Since $\mathbb{E}[\xi_t]=1$, the taste shock leaves the \emph{unconditional} ranking of keep versus adjust unchanged on average, but for any given $(m_t,n_t)$ there is now positive probability of adjusting even where $v_t^{\text{keep}} > v_t^{\text{adj}}(\ell_t)$, and vice versa, because a sufficiently large or small realization of $\xi_t$ can overturn the ranking implied by the two conditional values alone.

Because $\xi_{t+1}$ is drawn afresh every period, the object that matters for today's continuation value is the taste-shock-integrated value function
\begin{align}
	\bar v_{t+1}(m_{t+1},n_{t+1}) \equiv \mathbb{E}_{\xi_{t+1}}\Big[v_{t+1}(m_{t+1},n_{t+1},\xi_{t+1})\Big]
	= \mathbb{E}_{\xi_{t+1}}\Big[\max\big\{ v_{t+1}^{\text{keep}}(m_{t+1},n_{t+1}),\; \xi_{t+1}\, v_{t+1}^{\text{adj}}(\ell_{t+1}) \big\}\Big], \label{eq:taste_barv}
\end{align}
which replaces $v_{t+1}(m_{t+1},n_{t+1})$ inside the Bellman equations for $v_t^{\text{keep}}$ and $v_t^{\text{adj}}$. As $\sigma_\xi \to 0$, $\xi_t \to 1$ almost surely and the model collapses to the hard-maximum baseline of the previous subsection. Because $\bar v_{t+1}$ averages a kinked function over the taste-shock distribution, it is smoother than $v_{t+1}$ itself, and larger $\sigma_\xi$ smooths $\bar v_{t+1}$ (and the marginal-value functions derived from it) further, mitigating the propagation of kinks through the backward recursion. The taste shock does not alter the household's budget set, so the no-borrowing constraint $a_t \ge 0$ is unaffected. I derive the envelope conditions, the modified Euler equations, and the EGM algorithm with taste shocks in appendix~\ref{subsec:appdx_taste_basic}.

\subsection{Extension: Search Frictions}\label{subsec:ModelBasicSearch}

This subsection instead reintroduces the directed-search mechanism of the full model into the basic environment, without taste shocks. Rather than resolving the keep-versus-adjust margin with a hard maximum, the household chooses a search intensity $\varepsilon_t \in [0,1]$ that determines the weight placed on the adjuster's conditional value in the current period. Exerting search effort is costly, with cost function
\begin{align}
	\sigma(\varepsilon_t) = S\big[\varepsilon_t + (1-\varepsilon_t)\log(1-\varepsilon_t)\big], \label{eq:search_cost_basic}
\end{align}
where $S>0$ governs the degree of friction in the market for durables, following \cite{Khan2019-ew}. The beginning-of-period problem becomes
\begin{align}
	v_t(m_t,n_t) = \max_{\varepsilon_t \in [0,1]} \Big[(1-\varepsilon_t)\, v_t^{\text{keep}}(m_t,n_t) + \varepsilon_t\, v_t^{\text{adj}}(\ell_t) - \sigma(\varepsilon_t)\Big], \label{eq:search_discrete_basic}
\end{align}
where $v_t^{\text{keep}}$ and $v_t^{\text{adj}}$ retain exactly the same Bellman definitions as in the baseline model, except that the continuation value $v_{t+1}(m_{t+1},n_{t+1})$ appearing inside them is now itself given recursively by equation~\eqref{eq:search_discrete_basic} one period ahead, rather than by the hard maximum of the baseline model. All other features of the basic environment---the single interest rate $\mathcal{R}$, the absence of maintenance costs, and the no-borrowing constraint $a_t \ge 0$---are unchanged.

Search effort is chosen before the discrete outcome is known: $\varepsilon_t$ can be interpreted as the probability with which the household successfully locates a suitable durable this period, so the combined value function~\eqref{eq:search_discrete_basic} is a probability-weighted average of the two conditional values, net of the search cost. As $S \to 0$, search becomes costless, the household sets $\varepsilon_t \in \{0,1\}$, and the model collapses to the hard-maximum baseline; as $S \to \infty$, search becomes prohibitively costly, $\varepsilon_t^\ast \to 0$, and the household never adjusts. I derive the optimal search-effort rule, the smooth combined marginal-value functions, and the associated EGM steps in appendix~\ref{subsec:appdx_search_basic}.

\subsection{Extension: Stochastic Durable Prices}\label{subsec:ModelBasicPrice}

This subsection introduces a third source of variation, orthogonal to the taste shocks of Section~\ref{subsec:ModelBasicTaste} and the search frictions of Section~\ref{subsec:ModelBasicSearch}: an idiosyncratic, i.i.d.\ shock to the relative price at which the durable good can be purchased in the period of adjustment. Neither taste shocks nor search frictions are present in this subsection, so the discrete choice between keeping and adjusting again resolves through a hard maximum, but now the adjuster's problem---and hence the comparison itself---is contingent on the realized price.

At the start of period $t$, after the transitory income shock $\theta_t$ has materialized, the household observes an i.i.d.\ relative price $p_t$ at which new units of the durable good can be bought this period. The price is log-normally distributed with mean one,
\begin{align}
	\log p_t \sim \mathcal{N}\Big(-\frac{\sigma_p^2}{2},\, \sigma_p^2\Big), \qquad \mathbb{E}[p_t] = 1, \label{eq:price_dist}
\end{align}
independent of $\theta_t$ and independent across periods. Only the \emph{purchase} of new durables is exposed to this price risk: households liquidate their existing durable stock exactly as in the baseline model, at the fixed proportional recovery rate $1-\tau$, so the liquidity available upon adjustment, $\ell_t = m_t + (1-\tau)n_t$, is unaffected by $p_t$. This is consistent with an interpretation in which trading in a used durable is governed by a separate, deterministic adjustment cost, while the price of a new unit on the market this period fluctuates, e.g.\ a household trading in a used car against the fluctuating price of a new one.

The keeper's problem is entirely unaffected by $p_t$, since no purchase occurs, so $v_t^{\text{keep}}(m_t,n_t)$ retains its baseline Bellman definition. The adjuster's budget constraint, however, now reads
\begin{align}
	a_t = \ell_t - c_t - p_t\, d_t \ge 0, \label{eq:price_budget}
\end{align}
in place of $a_t=\ell_t-c_t-d_t$, so the conditional value function becomes a function of both $\ell_t$ and $p_t$,
\begin{align}
	v_t^{\text{adj}}(\ell_t,p_t) &= \max_{c_t,d_t} \ u(c_t,d_t) + \beta \mathbb{E}_t \big[ v_{t+1}(m_{t+1},n_{t+1}) \big],\\
	& \text{s.t. equation~\eqref{eq:price_budget},}\\
	& \qquad m_{t+1} = \mathcal{R} a_t + \theta_{t+1}, \qquad n_{t+1} = (1-\delta)\,d_t.
\end{align}
Everything else about the household's environment---the interest rate $\mathcal{R}$, the depreciation rate $\delta$, and the no-borrowing constraint $a_t \ge 0$---is unchanged, and next period's durable stock $n_{t+1}$ continues to be measured in physical units, so $p_t$ affects only the current-period cost of acquiring $d_t$, not its future service flow beyond ordinary depreciation.

Because $p_t$ is observed before the discrete choice is made, the household's beginning-of-period problem is the price-contingent hard maximum
\begin{align}
	v_t(m_t,n_t,p_t) = \max\Big\{ v_t^{\text{keep}}(m_t,n_t),\; v_t^{\text{adj}}(\ell_t,p_t)\Big\}. \label{eq:price_discrete}
\end{align}
A high realization of $p_t$ raises the marginal cost of adjusting and, other things equal, makes the household less likely to adjust and, conditional on adjusting, to purchase a smaller quantity of the durable; a low realization has the opposite effect. Because $\mathbb{E}[p_t]=1$, this leaves the \emph{average} cost of adjusting unchanged relative to the baseline model, but it introduces state-contingent variation in both the adjustment margin and the size of the resulting durable purchase.

As with the taste shock, the object that matters for the previous period's continuation value is the price-integrated value function
\begin{align}
	\bar v_t(m_t,n_t) \equiv \mathbb{E}_{p_t}\Big[v_t(m_t,n_t,p_t)\Big]
	= \mathbb{E}_{p_t}\Big[\max\big\{ v_t^{\text{keep}}(m_t,n_t),\; v_t^{\text{adj}}(\ell_t,p_t) \big\}\Big], \label{eq:price_barv}
\end{align}
which replaces $v_t(m_t,n_t)$ inside the Bellman equations for $v_{t-1}^{\text{keep}}$ and $v_{t-1}^{\text{adj}}$. Unlike the taste shock, which left the adjuster's optimal continuous choices unchanged and only reweighted the discrete comparison, the price shock changes the adjuster's optimal consumption-durable split itself, so $v_t^{\text{adj}}(\ell_t,p_t)$ must be solved separately for every relevant realization of $p_t$; the resulting policy functions $c_t^{\text{adj}}(\ell_t,p_t)$ and $d_t^{\text{adj}}(\ell_t,p_t)$ are likewise price-contingent, and the model's reported (ex-ante) nondurable-consumption and durable-purchase functions are price-weighted averages of these state-contingent choices, in the same spirit as $\bar v_t$. As $\sigma_p \to 0$, $p_t \to 1$ almost surely and the model collapses to the hard-maximum baseline of Section~\ref{sec:ModelBasic}. I derive the modified first-order conditions, the envelope conditions, the terminal-period solution, and the EGM algorithm with stochastic durable prices in appendix~\ref{subsec:appdx_price_basic}.

\subsection{Extension: Search, Taste Shocks, and Stochastic Prices}\label{subsec:ModelBasicFull}

This subsection nests the three extensions of Sections~\ref{subsec:ModelBasicTaste}--\ref{subsec:ModelBasicPrice} into a single environment. The household simultaneously faces (i)~an idiosyncratic taste shock $\xi_t$ that scales the value of adjusting, (ii)~an idiosyncratic relative price $p_t$ at which new durables can be purchased, and (iii)~a directed-search decision $\varepsilon_t\in[0,1]$ that determines the probability of successfully adjusting, at convex cost $\sigma(\varepsilon_t)$ from equation~\eqref{eq:search_cost_basic}. The three shocks are mutually independent and independent of income,
\begin{align}
	\log \xi_t &\sim \mathcal{N}\Big(-\frac{\sigma_\xi^2}{2},\,\sigma_\xi^2\Big), \qquad \mathbb{E}[\xi_t]=1, \label{eq:full_xi_dist}\\
	\log p_t &\sim \mathcal{N}\Big(-\frac{\sigma_p^2}{2},\,\sigma_p^2\Big), \qquad \mathbb{E}[p_t]=1, \label{eq:full_p_dist}
\end{align}
with $\xi_t\perp p_t\perp\theta_t$ and both i.i.d.\ across periods. As foreshadowed in Section~\ref{subsec:ModelBasicTaste}, once search is present the taste shock does not mechanically randomize the discrete outcome: it enters the household's optimal search-effort decision by scaling the adjuster's conditional value relative to the keeper's.

The within-period timeline is as follows. After $\theta_t$ has materialized and $(m_t,n_t)$ are known, the household observes the pair $(p_t,\xi_t)$. Liquidity upon adjustment remains $\ell_t=m_t+(1-\tau)n_t$, independent of both shocks. The keeper's problem is unaffected by either shock, so $v_t^{\text{keep}}(m_t,n_t)$ retains its baseline Bellman definition. The adjuster's budget constraint is the price-contingent constraint~\eqref{eq:price_budget}, so the adjuster's conditional value $v_t^{\text{adj}}(\ell_t,p_t)$ and its continuous policies $c_t^{\text{adj}}(\ell_t,p_t)$, $d_t^{\text{adj}}(\ell_t,p_t)$ are exactly as in Section~\ref{subsec:ModelBasicPrice}: $\xi_t$ does not enter the adjuster's inner maximization, while $p_t$ does. Given these two conditional values, the household chooses search intensity to maximize the taste-scaled, price-contingent objective
\begin{align}
	v_t(m_t,n_t,p_t,\xi_t)
	&= \max_{\varepsilon_t\in[0,1]}
	\Big[(1-\varepsilon_t)\, v_t^{\text{keep}}(m_t,n_t)
	+ \varepsilon_t\, \xi_t\, v_t^{\text{adj}}(\ell_t,p_t)
	- \sigma(\varepsilon_t)\Big]. \label{eq:full_discrete}
\end{align}
Equation~\eqref{eq:full_discrete} nests each of the three standalone extensions as a special case: with $S\to 0$ and $\sigma_p\to 0$ it collapses to the taste-shock hard maximum~\eqref{eq:taste_discrete}; with $\sigma_\xi\to 0$ and $\sigma_p\to 0$ it collapses to the search problem~\eqref{eq:search_discrete_basic}; with $S\to 0$ and $\sigma_\xi\to 0$ it collapses to the price-contingent hard maximum~\eqref{eq:price_discrete}; and with all three limits simultaneous, $S\to 0$, $\sigma_\xi\to 0$, $\sigma_p\to 0$, it collapses to the hard-maximum baseline of Section~\ref{sec:ModelBasic}.

Because $(p_{t+1},\xi_{t+1})$ are drawn afresh every period, the object that enters today's Euler equations is the jointly integrated value function
\begin{align}
	\bar v_t(m_t,n_t)
	&\equiv \mathbb{E}_{p_t,\xi_t}\big[v_t(m_t,n_t,p_t,\xi_t)\big], \label{eq:full_barv}
\end{align}
which replaces $v_t(m_t,n_t)$ inside the Bellman equations for $v_{t-1}^{\text{keep}}$ and $v_{t-1}^{\text{adj}}$. The three frictions play complementary roles: $p_t$ shifts the adjuster's continuous problem and therefore the level of $v_t^{\text{adj}}$; $\xi_t$ reweights that level in the search objective without changing the adjuster's continuous choices; and $\varepsilon_t^\ast$ translates the resulting gap $ \xi_t v_t^{\text{adj}}-v_t^{\text{keep}}$ into a smooth probability of adjusting. I derive the optimal search-effort rule, the combined envelope conditions, the Euler equations, the terminal-period solution, and the EGM algorithm for the nested model in appendix~\ref{subsec:appdx_full_basic}.
